% !TEX program = pdflatex

%%%%%%%%%%%%%%%%%%%%
%affineConway.tex
%%%%%%%%%%%%%%%%%%%%%

\pdfoutput=1 %When you use pdflatex
\documentclass{amsart} %When you use pdflatex
\renewcommand{\baselinestretch}{1.2}

\usepackage{tikz}

\usepackage{amssymb}
\usepackage{amsmath}
\usepackage{url}
\usepackage{bm}
\usepackage{enumitem}
\usepackage{graphicx}

\usepackage{hyperref}
\hypersetup{
 colorlinks = true,
 linkcolor = blue,
 filecolor = blue,
 urlcolor = blue,
 citecolor = blue,
}
\hypersetup{unicode=true}

%\input mymacros.tex
%%%%%%%%%%%%%%%%%%%%%%%%%%%%%%%%%%%%%%%%%%%%
%
%   MY MACROS: start 2025 Aug 12
%
%%%%%%%%%%%%%%%%%%%%%%%%%%%%%%%%%%%%%%%%%%%%%

% erase: for temporarily removing text from output
\newcommand{\erase}[1]{}

% ----- Theorem-like environments -----
% Plain style: italic body, bold heading
\theoremstyle{plain}
\newtheorem{theorem}{Theorem}[section]       % Main theorem environment
\newtheorem{lemma}[theorem]{Lemma}           % Lemma, shares numbering with theorem
\newtheorem{proposition}[theorem]{Proposition} % Proposition, shares numbering
\newtheorem{corollary}[theorem]{Corollary}   % Corollary, shares numbering

% Definition style: upright body, bold heading
\theoremstyle{definition}
\newtheorem{definition}[theorem]{Definition} % Definition, shares numbering
\newtheorem{example}[theorem]{Example}       % Example
\newtheorem{algorithm}[theorem]{Algorithm}   % Algorithm description
\newtheorem{claim}[theorem]{Claim}           % Claim
\newtheorem{construction}[theorem]{Construction} % Construction

% Remark style: upright body, italic heading
\theoremstyle{remark}
\newtheorem{remark}[theorem]{Remark}         % Remark

% Number equations, tables, and figures by section
\numberwithin{equation}{section}
\numberwithin{table}{section}
\numberwithin{figure}{section}

% Change QED symbol for end of proof
\newcommand{\qedbox}{\hfill {$\Box$}}

% ----- FONT SHORTCUTS -----
% Blackboard bold (mathbb)
\newcommand{\bbA}{\mathord{\mathbb{A}}}
\newcommand{\BB}{\mathord{\mathbb{B}}}
\newcommand{\CC}{\mathord{\mathbb{C}}}
\newcommand{\DD}{\mathord{\mathbb{D}}}
\newcommand{\EE}{\mathord{\mathbb{E}}}
\newcommand{\FF}{\mathord{\mathbb{F}}}
\newcommand{\GG}{\mathord{\mathbb{G}}}
\newcommand{\HH}{\mathord{\mathbb{H}}}
\newcommand{\II}{\mathord{\mathbb{I}}}
\newcommand{\JJ}{\mathord{\mathbb{J}}}
\newcommand{\KK}{\mathord{\mathbb{K}}}
\newcommand{\LL}{\mathord{\mathbb{L}}}
\newcommand{\MM}{\mathord{\mathbb{M}}}
\newcommand{\NN}{\mathord{\mathbb{N}}}
\newcommand{\OO}{\mathord{\mathbb{O}}}
\newcommand{\PP}{\mathord{\mathbb{P}}}
\newcommand{\QQ}{\mathord{\mathbb{Q}}}
\newcommand{\RR}{\mathord{\mathbb{R}}}
\newcommand{\bbSS}{\mathord{\mathbb{S}}}
\newcommand{\TT}{\mathord{\mathbb{T}}}
\newcommand{\UU}{\mathord{\mathbb{U}}}
\newcommand{\VV}{\mathord{\mathbb{V}}}
\newcommand{\WW}{\mathord{\mathbb{W}}}
\newcommand{\XX}{\mathord{\mathbb{X}}}
\newcommand{\YY}{\mathord{\mathbb{Y}}}
\newcommand{\ZZ}{\mathord{\mathbb{Z}}}

% Calligraphic letters (mathcal)
\newcommand{\AAA}{\mathord{\mathcal{A}}}
\newcommand{\BBB}{\mathord{\mathcal{B}}}
\newcommand{\CCC}{\mathord{\mathcal{C}}}
\newcommand{\DDD}{\mathord{\mathcal{D}}}
\newcommand{\EEE}{\mathord{\mathcal{E}}}
\newcommand{\FFF}{\mathord{\mathcal{F}}}
\newcommand{\GGG}{\mathord{\mathcal{G}}}
\newcommand{\HHH}{\mathord{\mathcal{H}}}
\newcommand{\III}{\mathord{\mathcal{I}}}
\newcommand{\JJJ}{\mathord{\mathcal{J}}}
\newcommand{\KKK}{\mathord{\mathcal{K}}}
\newcommand{\LLL}{\mathord{\mathcal{L}}}
\newcommand{\MMM}{\mathord{\mathcal{M}}}
\newcommand{\NNN}{\mathord{\mathcal{N}}}
\newcommand{\OOO}{\mathord{\mathcal{O}}}
\newcommand{\PPP}{\mathord{\mathcal{P}}}
\newcommand{\QQQ}{\mathord{\mathcal{Q}}}
\newcommand{\RRR}{\mathord{\mathcal{R}}}
\newcommand{\SSS}{\mathord{\mathcal{S}}}
\newcommand{\TTT}{\mathord{\mathcal{T}}}
\newcommand{\UUU}{\mathord{\mathcal{U}}}
\newcommand{\VVV}{\mathord{\mathcal{V}}}
\newcommand{\WWW}{\mathord{\mathcal{W}}}
\newcommand{\XXX}{\mathord{\mathcal{X}}}
\newcommand{\YYY}{\mathord{\mathcal{Y}}}
\newcommand{\ZZZ}{\mathord{\mathcal{Z}}}


% Fraktur letters (mathfrak)
\newcommand{\SSSS}{\mathord{\mathfrak{S}}}    % Symmetric group symbol
\newcommand{\AAAA}{\mathord{\mathfrak{A}}}    % Alternating group symbol

%For vectors
\newcommand{\vect}[1]{\boldsymbol{#1}}

% ----- ARROWS -----

%ARROWS

\newcommand{\maprightsp}[1]{\xrightarrow{\smash{#1}}}
\newcommand{\mapleftsp}[1]{\xleftarrow{\smash{#1}}}
\newcommand{\maprightsb}[1]{\xrightarrow[\smash{#1}]{}}
\newcommand{\maplefts}[1]{\xleftarrow[\smash{#1}]{}}

\newcommand{\inj}{\hookrightarrow}
\providecommand{\twoheadrightarrow}{\mathrel{\rightarrow\!\!\!\!\rightarrow}}
\newcommand{\surj}{\twoheadrightarrow}

\newcommand{\isom}{\xrightarrow{\raise -3pt \hbox{\scriptsize $\sim$} }}


% Vertical arrows
\newcommand{\mapdown}{\mathrel{\downarrow}}
\newcommand{\mapdownleft}[1]{\mathrel{\llap{\scriptstyle#1\;}\downarrow}}


% ----- SET NOTATION -----
\newcommand{\set}[2]{\{\,{#1}\mid {#2} \,\}}           % { x | condition }
\newcommand{\bigset}[2]{\left\{\; {#1} \; \left\vert \; {#2} \;  \right.\right \}}
\newcommand{\bigsetR}[2]{\left\{\; \left. {#1} \; \right\vert \; {#2} \;  \right \}}

\newcommand{\angs}[1]{\langle {#1}  \rangle}           % < x >


% ----- COMMON SYMBOLS -----
\newcommand{\wt}{\widetilde}
\newcommand{\ol}{\overline}
\newcommand{\tensor}{\otimes}
\newcommand{\inv}{^{-1}}                 % superscript -1
\newcommand{\dual}{^{\vee}}               % dual vector space
\newcommand{\comp}{^{\wedge}}             % complement or wedge power

% Primes and related notations
\newcommand{\pprime}{\prime\prime}
\newcommand{\sprime}{^{\prime}}
\newcommand{\sprimeinv}{^{\prime-1}}
\newcommand{\spar}[1]{^{(#1)}}
\newcommand{\spprime}{^{\prime\prime}}
\newcommand{\spcirc}{^{\mathord{\circ}}}
\newcommand{\sptimes}{^{\times}}
\newcommand{\sperp}{^{\perp}}

% Group theory and algebra notation
\newcommand{\fiberproduct}{\Box}            % fibre product symbol
\newcommand{\semidirectproduct}{\rtimes}    % semidirect product
\newcommand{\semidirectproductrev}{\ltimes} 

% Boundary and derivative
\newcommand{\bdr}{\partial\,}
\newcommand{\der}{\partial}

% Math operators (upright text in math mode)
\DeclareMathOperator{\Aut}{Aut}
\DeclareMathOperator{\Image}{Im}
\DeclareMathOperator{\Sing}{Sing}
\DeclareMathOperator{\Spec}{Spec}
\DeclareMathOperator{\Stab}{Stab}
\DeclareMathOperator{\Ker}{Ker}
\DeclareMathOperator{\Coker}{Coker}
\DeclareMathOperator{\Gal}{Gal}
\DeclareMathOperator{\Hom}{Hom}
\DeclareMathOperator{\rank}{rank}
\DeclareMathOperator{\pr}{pr}
\DeclareMathOperator{\disc}{disc}
\DeclareMathOperator{\sgn}{sgn}
\DeclareMathOperator{\characteristic}{char}
\DeclareMathOperator{\length}{leng} % 'leng' probably intended as 'length'

% Sheaf Hom
\newcommand{\sheafHom}{\mathord{\mathcal {Hom}}}

% Classical groups
\newcommand{\GL}{\mathord{\mathrm{GL}}}
\newcommand{\SL}{\mathord{\mathrm{SL}}}
\newcommand{\PGU}{\mathord{\mathrm{PGU}}}
\newcommand{\PSU}{\mathord{\mathrm{PSU}}}
\newcommand{\PGL}{\mathord{\mathrm{PGL}}}
\newcommand{\PSL}{\mathord{\mathrm{PSL}}}
\newcommand{\OG}{\mathord{\mathrm{O}}}
\newcommand{\UG}{\mathord{\mathrm{U}}}
\newcommand{\SO}{\mathord{\mathrm{SO}}}
\newcommand{\SU}{\mathord{\mathrm{SU}}}

% Identity map
\newcommand{\id}{\mathord{\mathrm{id}}}
\newcommand{\Id}{\mathord{\mathrm{Id}}}

% Common notation shortcuts
\newcommand{\Pt}{\PP^2}                      % Projective plane
\newcommand{\Qbar}{\overline{\QQ}}           % Algebraic closure of Q
\newcommand{\pione}{\pi_1}  % Fundamental group

% Inner product or pairing
\newcommand{\intf}[1]{\langle #1 \rangle}

% Transpose symbol
\newcommand{\transpose}{{}^{\mathrm{T}}}

% Vertical spacing helpers (artificially increase the vertical spacing )
\newcommand{\mystruth}[1]{\phantom{\llap{\vrule height #1}}}
\newcommand{\mystrutd}[1]{\phantom{\llap{\vrule depth #1}}}
\newcommand{\mystruthd}[2]{\phantom{\llap{\vrule  height #1 depth #2}}}



%%%%%%%%%%%%%%%%%%%%%%%%%%%%%%%%%%%%%%%%%%%%
%
%   END OF MY MACROS
%
%%%%%%%%%%%%%%%%%%%%%%%%%%%%%%%%%%%%%%%%%%%%%
%\input mymacrosaffineConway.tex
%%%%%%%%%%%%%%%%%%%%%%%%%%%%%%%%%%%%%%%%%%%%
%
%   MY MACROS: start 
%
%%%%%%%%%%%%%%%%%%%%%%%%%%%%%%%%%%%%%%%%%%%%%

\newcommand{\dotz}{\cdot 0}
\newcommand{\dotinf}{\cdot \infty}
\newcommand{\Invols}{\III}
\newcommand{\Cen}{\mathrm{Cen}}
\newcommand{\Leech}{\Lambda}
\newcommand{\Golay}{\CCC_{24}}
\newcommand{\pig}{\pi_{g}}
\newcommand{\varpigt}{\varpi_{(g, \tau)}}
\newcommand{\vep}{\varepsilon}
\newcommand{\PowOmega}{2^{\Omega}}
\newcommand{\weyl}{\vect{w}}

\newcommand{\IIlat}{\mathrm{II}}
\newcommand{\typeILat}{\mathrm{I}}
\newcommand{\typeIILat}{\mathrm{II}}
\newcommand{\Lts}{L_{26}}
\newcommand{\PPPts}{\PPP_{26}}
\newcommand{\FFFts}{{\FFF}_{26}}
\newcommand{\enrinvol}{\epsilon}

\newcommand{\Sigmatwlv}{\varSigma_{12}}
\newcommand{\Laminated}{\varLambda}
\newcommand{\BW}{\mathrm{BW}}

\newcommand{\tname}[1]{$\mathtt{#1}$}


\newcommand{\Lamsxtn}{\Laminated_{16}}
\newcommand{\barLeech}{\overline{\Leech}}
\newcommand{\compword}[1]{\overline{#1}}

\newcommand{\intfLeech}[1]{\intf{#1}_{\Leech}}
\newcommand{\intfL}[1]{\intf{#1}_{\hskip -1pt L}}
\newcommand{\intfM}[1]{\intf{#1}_{\hskip -1pt M}}
\newcommand{\intfSigma}[1]{\intf{#1}_{\hskip -1pt {\varSigma}}}

\newcommand{\tilgamma}{\tilde{\gamma}}
\newcommand{\tilg}{\tilde{g}}

\newcommand{\LeechRoot}[1]{r(#1)}

\newcommand{\AOG}{\mathord{\mathrm{AO}}}

\newcommand{\ConCham}{\mathord{C}}
\newcommand{\ConChamz}{\ConCham(\weyl_0)}


\newcommand{\walls}{\WWW}
\newcommand{\tilwalls}{\widetilde{\WWW}}
\newcommand{\tilXi}{\widetilde{\Xi}}
\newcommand{\Xitau}{\Xi_{\tau}}
\newcommand{\tilXitau}{\tilXi_{\tau}}

\newcommand{\Lplus}{L_{+}}
\newcommand{\Lminus}{L_{-}}
\newcommand{\qplus}{q_{+}}
\newcommand{\qminus}{q_{-}}
\newcommand{\gplus}{g_{+}}
\newcommand{\gminus}{g_{-}}
\newcommand{\Leechplus}{\Leech_{+}}
\newcommand{\Leechminus}{\Leech_{-}}
\newcommand{\etaplus}{{\eta}_{+}}
\newcommand{\etaminus}{{\eta}_{-}}
\newcommand{\lambdaplus}{{\lambda}_{+}}
\newcommand{\lambdaminus}{{\lambda}_{-}}
\newcommand{\FFFplus}{\FFF_{+}}
\newcommand{\prplus}{\pr_{+}}

\newcommand{\PPPplus}{\PPP_{+}}
\newcommand{\vGam}{\varGamma}
\newcommand{\tilvGam}{\widetilde{\varGamma}}
\newcommand{\vSig}{\varSigma}

\newcommand{\prM}{\pr_M}
\newcommand{\prS}{\pr_S}


\newcommand{\prLplus}{\pr^{L}_{+}}
\newcommand{\prLminus}{\pr^{L}_{-}}
\newcommand{\prLeechplus}{\pr^{\Leech}_{+}}
\newcommand{\prLeechminus}{\pr^{\Leech}_{-}}

\newcommand{\barSigma}{\overline{\Sigma}}

%\newcommand{\tauI}{\tau_{\,\typeILat}}
%\newcommand{\tauII}{\tau_{\,\typeIILat}}
\newcommand{\tauI}{\tau_{120}}
\newcommand{\tauII}{\tau_{135}}
%%%%%%%%%%%%%%%%%%%%%%%%%%%%%%%%%%%%%%%%%%%%
%
%   END OF MY MACROS
%
%%%%%%%%%%%%%%%%%%%%%%%%%%%%%%%%%%%%%%%%%%%%%

\newcommand{\XC}{X_C}
\newcommand{\diag}{\mathrm{diag}}
%%%%%%%%%%%%%%%%%%%%%%%%%%%%
%
\newcommand{\ps}{B}
\newcommand{\con}{C}
\newcommand{\wtcon}{\widetilde{\con}}
\newcommand{\pips}{\pi_{\ps}}
\newcommand{\Xps}{X_{\ps}}
\newcommand{\col}{\Gamma }
\newcommand{\gra}{w}
\newcommand{\varpips}{\varpi_{\ps}}
\newcommand{\Pic}{\mathrm{Pic}}
\newcommand{\bPP}{\mathbf{P}}
\newcommand{\Sw}{\mathrm{Sw}}
\newcommand{\Wk}{\mathcal{W}_k}
%\newcommand{\abs}[1]{ {#1}^{\mathrm{abs}}}
\newcommand{\abs}[1]{\mathrm{ab}(#1)}
\newcommand{\ab}{\mathrm{ab}}
\newcommand{\lex}{\mathrm{lex}}
\newcommand{\smallerlex}{<_{\mathrm{lex}}}
%
\begin{document}
%
\title{Remarks on  six-tangential conics to smooth plane sextics} 
\author{Ichiro Shimada}
\address{Department of Mathematics,
Graduate School of Science,
Hiroshima University,
1-3-1 Kagamiyama,
Higashi-Hiroshima,
739-8526 JAPAN}
\email{ichiro-shimada@hiroshima-u.ac.jp}
%
%超平面配置に関連する離散構造の拡張、深化とその応用
%23H00081
%代数幾何学の計算機による研究の新展開
%23K20209
%格子、保型形式とK3曲面、エンリケス曲面の研究
%20H00112
%

%
\begin{abstract}
\end{abstract}
%

\keywords{}
\maketitle

Let $k$ be a positive integer.
We consider the complete graph with vertex set $V_k=\{1, \dots, k\}$ and edge set 
\[
E_k:=\set {[i, j]}{i, j\in V_k, i<j}.
\]
Let $m$ be a positive integer, and we put 
\[
M:=\{-m,-m+1, \dots, m\}.
\]
Let $\Wk$ be the set of weights $w \colon E_k \to M$, 
that is, $\Wk$ is the set of weighted complete graphs on $V_k$
with weights in $M$.
We let  the symmetric group $\SSSS_k$ act on $V_k$ from the right.
Then $\SSS_k$ acts on $\Wk$ from the left.
We  define $\Sw_k$ to be the multiplicative group of all mappings
$s \colon V_k \to \{\pm 1\}$, 
which acts on $\Wk$ by switching, 
that is, 
\[
(sw) ([i, j])= s(i) s(j)  w([i, j])
\]
for $s\in \Sw_k$ and $[i, j]\in E_k$.
Then the semidirect product
\[
G_k := \Sw_k \rtimes \SSSS_k
\]
acts on $\Wk$ from the left.
The purpose of this section is to describe an efficient algorithm
for obtaining the list of 
orbits of the action of $G_k$ on $\Wk$.

First, we introduce a total order on the set $E_k$ of edges by
\[
[i,j] < [i\sprime,j\sprime]
\quad\Longleftrightarrow\quad
j < j\sprime
\quad\text{or}\quad
(j = j\sprime \text{ and } i < i\sprime).
\]
Then we express elements $w$ of $\Wk$ by arranging
the values $w([i,j])$ in this order:
\[
w = \bigl(w([1,2]), w([1,3]), w([2,3]), w([1,4]), \ldots \bigr).
\]
Let $\abs{w}$ denote the vector obtained from $w$ by taking
the absolute value of each component.
We put $\abs{M} := \{0,1,\ldots,m\}$.
Then we have a mapping
%
\begin{equation}\label{eq:absw}
\abs \colon \Wk \longrightarrow \abs{\Wk} := (\abs{M})^{E_k},
\qquad
w \longmapsto \abs{w}.
\end{equation}
%
\begin{definition}
Suppose that a group $G$ acts on a set $X$,
and that $X$ is equipped with 
a total order
 $<_{X}$.
An element $x\in X$ is \emph{$(G, <_{X})$-minimal} 
if $x$ is minimal 
with respect to $<_{X}$ in the orbit $G\cdot x$.
\end{definition}
%
Notice that $\SSSS_k$ acts on $\abs{\Wk}$, and the mapping~\eqref{eq:absw}
is $\SSSS_k$-equivariant.
Moreover, the action of $\Sw_k$ on $\Wk$ preserves
each fiber of~\eqref{eq:absw}.
We then define a total order $<_{\WWW}$ on $\Wk$ by
\[
w <_{\WWW} w\sprime
\quad\Longleftrightarrow\quad
\abs{w} <_{\mathrm{lex}} \abs{w\sprime}
\quad\text{or}\quad
\bigl(\abs{w} = \abs{w\sprime} \;\;\text{ and }\;\; w <_{\mathrm{lex}} w\sprime\bigr).
\]
We obtain the list of $G_k$-orbits in $\Wk$
by enumerating all $(G_k, <_{\WWW} )$-minimal elements in $\Wk$.
\par
We define the notion of  
 \emph{$G_k$-quasi-minimality}
by the following:
%$w$ is $G_k$-minimal if and only if $\abs{w}$ is minimal
%in its $G_k$-orbit $G_k \abs{w}$ in $\abs{\Wk}$, and $w$ is minimal
%in its $\Sw_k$-orbit  $\Sw_k w$.
\[
\text{$w$ is $G_k\text{-quasi-minimal}$}
\quad\Longleftrightarrow\quad
\parbox{7cm}{
$\abs{w}$ is$(\SSSS_k, <_{\mathrm{lex}})$-minimal, and  \\ 
$w$  is $(\Sw_k, <_{\mathrm{lex}})$-minimal   in  $\ab\inv(\ab(w))$.}
\]
%
It is obvious that, if  $w \in \Wk$ is $(G_k, <_{\WWW} )$-minimal then  $w \in \Wk$ is $G_k$-quasi-minimal
while the opposite implication does not necessarily hold.
We first make the list of $G_k$quasi--minimal elements,
and then we choose from them $(G_k, <_{\WWW} )$-minimal elements.
\par
To make the list of $G_k$-quasi-minimal elements,
we use the induction on $k$.
Suppose that $k\sprime < k$.
By taking the prefix of length $k\sprime(k\sprime-1)/2$ of each $u \in \abs{\Wk}$,
we obtain a mapping
\[
\abs{\Wk} \longrightarrow \abs{\WWW_{k\sprime}},
\qquad
u \longmapsto u|_{k\sprime}.
\]
It is obvious that, if $u\in \abs{\Wk}$ is $(\SSSS_k, <_{\lex})$-minimal,
then $u|_{k\sprime}$ must be  $(\SSSS_{k\sprime}, <_{\lex})$-minimal.
 \par
 For an $(\SSSS_k, <_{\lex})$-minimal element $u\in \abs{\Wk}$, 
 \[
 \ab\inv (u)  \longrightarrow \ab\inv(u|_{k\sprime}),
 \qquad
w \longmapsto w|_{k\sprime}.
 \]
 It is obvious that, if $w\in \ \ab\inv (u)$ is $(\Sw_k, <_{\lex})$-minimal,
then $w|_{k\sprime}$ must be  $(\Sw_{k\sprime}, <_{\lex})$-minimal.
More strongly, 
$w\in \ \ab\inv (u)$ is $(\Sw_k, <_{\lex})$-minimal
only if 
$w|_{k\sprime}$ is $(\Sw_{k\sprime}, <_{\lex})$-minimal
and the first non-zero element of the suffix $w-w|_{k\sprime}$ is negative.
(The opposite implication does not necessarily hold.)
We say that $w$ is $(\SSSS_k, <_{\lex})$-quasi-minimal 
if, for any $k\sprime<k$,
 $w|_{k\sprime}$ is  $(\SSSS_k, <_{\lex})$-quasi-minimal 
and the first non-zero element of the suffix $w-w|_{k\sprime}$ is negative.
\par
We enumerate all  $(\SSSS_k, <_{\lex})$-minimal elements 
and then, for each  $(\SSSS_k, <_{\lex})$-minimal element $u\in \abs{\Wk}$,
we enumerate all  $(\Sw_k, <_{\lex})$-quasi-minimal  elements in $\ab\inv (u)$.
%
%
\end{document}


